% ECONOMICS_PROBLEM: {"id": "C23-372", "title": "Synthetic-control inference with few treated units", "jel_code": "C23", "source_id": "OP-0106"}
% !TeX program = xelatex
\documentclass[11pt,a4paper]{article}
\usepackage[margin=23mm]{geometry}
\usepackage{fontspec}
\setmainfont{Latin Modern Roman}
\usepackage{amsmath,amssymb,mathtools}
\usepackage{xcolor,enumitem,booktabs,longtable,array,xurl,hyperref}
\definecolor{muted}{HTML}{53636C}
\hypersetup{colorlinks=true,urlcolor=blue,linkcolor=blue}
\setlength{\parindent}{0pt}
\setlength{\parskip}{5pt}
\newcommand{\R}{\mathbb R}
\newcommand{\E}{\mathbb E}
\newcommand{\N}{\mathbb N}
\newcommand{\ind}{\mathbf 1}
\newcommand{\Var}{\operatorname{Var}}
\newcommand{\Cov}{\operatorname{Cov}}
\newcommand{\supp}{\operatorname{supp}}
\DeclareMathOperator*{\argmin}{arg\,min}
\DeclareMathOperator*{\argmax}{arg\,max}
\newcommand{\entrymeta}[3]{{\sffamily\small\color{muted}\textbf{\texttt{#1}}\quad|\quad #2\par #3\par}}
\newcommand{\Problemref}[1]{\texttt{#1}}
\newcommand{\JELref}[2]{\texttt{#2} (JEL \texttt{#1})}
\begin{document}
\section*{Synthetic-control inference with few treated units}
\textbf{Problem ID:} \texttt{C23-372}\quad \textbf{JEL:} \texttt{C23}\par
% BEGIN ECONOMICS STATEMENT
\label{problem:OP-0106}
\label{atlas:prob:E6}

\noindent\textbf{Original identifier:} E6 (atlas item 106). \textbf{Field:} Econometrics and causal inference. \textbf{Classification:} editorial research family with established restricted solutions; current open status not certified.

\subsection*{Definitions and assumptions}

Observe a panel $Y_{it}$ for units $i=1,\ldots,N+1$ and times $t=1,\ldots,T$, with unit $1$ treated only after $T_0<T$ and other units untreated. Let $Y_{it}(d)$ be outcomes under a specified treatment regime $d\in\{0,1\}$, with no anticipation and no spillovers to donors. A candidate untreated-outcome model is $Y_{it}(0)=a_i+b_t+\lambda_i^\top f_t+u_{it}$, where $\lambda_i,f_t\in\mathbb R^r$, $r$ is fixed or has stated growth, and $u_{it}$ follows a specified time-series law. The realized post-treatment contrast is $\theta=(T-T_0)^{-1}\sum_{t>T_0}[Y_{1t}(1)-Y_{1t}(0)]$. Observed outcomes equal $Y_{1t}(1)$ for $t>T_0$ and $Y_{it}(0)$ otherwise.

\subsection*{Formal research question}

For a declared stochastic-panel or random-assignment experiment, construct intervals for $\theta$ that include donor selection, counterfactual-model ambiguity and post-treatment prediction noise. Let $\mathcal P_{N,T}$ restrict factor strength, temporal dependence, donor contamination and an approximation error $\Delta_{N,T}$ for representing unit $1$ by donor weights $w_j\ge0$, $\sum_{j=2}^{N+1}w_j=1$. Seek
\[
\liminf_{N,T\to\infty}\inf_{P\in\mathcal P_{N,T}}
\Pr_P\{\theta\in C_{N,T}\}\ge1-\alpha,\qquad\alpha\in(0,1),
\]
with explicit relative growth of $T_0$, $T-T_0$ and $N$. Determine when diameter decreases, and lower bounds when an unspanned factor or a single unpredictable post-treatment shock cannot be learned. Also give finite-sample randomization guarantees if the treated unit was actually assigned by a known design.

\subsection*{Interpretation and scope}

The coverage probability must be tied to a real experiment; with a fixed treated country there is no automatic exchangeability of its treatment label with donor labels. Placebo ranks may be diagnostics without being exact randomization $p$-values. Good pretreatment fit does not prove that latent factor relationships persist after treatment. Inference for a conditional mean effect differs from predicting a realized untreated trajectory, which includes irreducible idiosyncratic noise. A growing donor pool cannot necessarily compensate for a fixed post-treatment horizon. Negative donor weights, extrapolation and matrix completion require distinct approximation assumptions.

\subsection*{Source audit and status}

[S1] and [S2] give established synthetic-control constructions and applications; [S3] develops synthetic difference-in-differences. Their assumptions and asymptotic regimes should not be converted into a universal one-treated-unit guarantee. The residual research question is scoped robustness to model choice and limited post-treatment information. This is an editorial formulation supported by the literature, with no certified claim that all its variants remain open in 2026.

\subsection*{References}

\begin{enumerate}[label={[\textnormal{S}\arabic*]},leftmargin=*]

\item Alberto Abadie, Alexis Diamond and Jens Hainmueller (2010). \emph{Synthetic Control Methods for Comparative Case Studies: Estimating the Effect of California’s Tobacco Control Program}. Journal of the American Statistical Association 105(490), 493–505. \url{https://doi.org/10.1198/jasa.2009.ap08746}. \textbf{Locator:} Publisher or institutional abstract and bibliographic record. \textbf{Support and limits:} Synthetic-control construction and comparative case study; placebo ranks need a defensible assignment/exchangeability interpretation.

\item Alberto Abadie, Alexis Diamond and Jens Hainmueller (2015). \emph{Comparative Politics and the Synthetic Control Method}. American Journal of Political Science 59(2), 495–510. \url{https://doi.org/10.1111/ajps.12116}. \textbf{Locator:} Publisher or institutional abstract and bibliographic record. \textbf{Support and limits:} Methodological comparative-case analysis; not universal small-sample uncertainty.

\item Dmitry Arkhangelsky, Susan Athey, David A. Hirshberg, Guido W. Imbens and Stefan Wager (2021). \emph{Synthetic Difference-in-Differences}. American Economic Review 111(12), 4088–4118. \url{https://doi.org/10.1257/aer.20190159}. \textbf{Locator:} Publisher or institutional abstract and bibliographic record. \textbf{Support and limits:} Panel estimator with asymptotic conditions; not an unrestricted one-treated-unit inference theorem.

\end{enumerate}

\subsection*{Common conventions: Source mathematical conventions}

Unless an entry states otherwise, agent and firm sets are finite. State and action spaces are measurable, strategies and outcome maps are measurable, probability laws are specified, and expectations exist for the targets under discussion. Infinite-horizon discount factors lie in $(0,1)$ and discounted objectives are finite. Norms, tolerances and complexity input sizes must be specified for computational comparisons. Constants or primitives that appear locally are inputs, not empirical facts asserted by this catalogue.

\textbf{Identification.} A statistical model $m\in\mathcal M$ implies an observable law $P_m$ and target $\theta(m)$. At law $P$, its identified set is
\[
\Theta(P;\mathcal M)=\{\theta(m):m\in\mathcal M,\ P_m=P\}.
\]
Point identification holds when this set is a singleton, or equivalently when $P_{m_1}=P_{m_2}$ implies $\theta(m_1)=\theta(m_2)$ for all compatible models. Sharp bounds mean bounds attained, or approached when the boundary is not attained, by compatible models. Writing this set is a definition, not a solution: the substantive task is to establish informative restrictions and derive or estimate the set. An unrestricted-confounding model may give only trivial bounds.

\textbf{Inference.} Identification, estimability and valid uncertainty quantification are separate questions. Fix $\alpha\in(0,1)$. A confidence procedure $C_n$ over a declared law class $\mathcal P$ has asymptotic uniform coverage at level $1-\alpha$ when
\[
\liminf_{n\to\infty}\inf_{P\in\mathcal P}\Pr_P\{\theta(P)\in C_n\}\geq1-\alpha.
\]
For set-identified targets the coverage event must be defined separately, for example coverage of the full identified set. If an entry uses identification-indexed models instead of uniquely indexed laws, the same coverage convention is applied over those models. Pointwise limits do not establish uniform validity, and asymptotic validity does not guarantee finite-sample precision. Sample sizes, dependence structures and weak-identification sequences are part of each application.

\textbf{Counterfactuals.} $Y_i(a)$ is a potential outcome under a specified intervention, including its timing, implementation and versions. It is not $Y_i$ conditional on voluntarily choosing $a$. A full intervention vector permits interference; shorthand scalar treatment notation does not silently impose no interference. Assignment, selection, attrition, anticipation and measurement must be specified. A claim about an instrument's local effect is not automatically a population policy effect.

\textbf{Economic equilibrium.} A counterfactual model includes best responses, constraints, feasibility, market clearing, state transitions and expectations consistent with the stated information. When several equilibria exist, either a declared selector is an input or the target is set-valued. Comparisons across policy regimes must use the stated selection convention. Reduced-form relationships are not presumed invariant after an equilibrium-changing intervention.

\textbf{Optimization and welfare.} Optimization is over the stated feasible set. If attainment is not established, $\sup$ or $\inf$ and $\varepsilon$-optimal choices replace an asserted maximizer or minimizer. For example, with value $V=\sup_{a\in\mathcal A}W(a)<\infty$, an $\varepsilon$-optimal set is $\{a\in\mathcal A:W(a)\geq V-\varepsilon\}$ for $\varepsilon>0$. Benefits, costs, risk and health or environmental outcomes can be aggregated only using declared units and conversion weights. Interpersonal utility weights, the treatment of fiscal transfers and foreign profits, discounting, and the behavioral welfare criterion are explicit inputs. A comparative-static derivative additionally requires existence and appropriate differentiability of the selected counterfactual.

\textbf{Sources.} Local labels [S1], [S2], and so forth restart in every question. Locators identify the relevant abstract, model, section or result at the level actually checked; a bibliographic or abstract-level verification is not represented as inspection of an entire proof. Working papers are distinguished from later publications and changing versions. Source summaries are paraphrased. All URL checks and editorial formulations refer to the audit date stated above.
% END ECONOMICS STATEMENT
\end{document}
